\chapter{Literature Review}
\label{ch:literature}

\section{Overview}

This chapter traces the development of neural text classification in order to establish where the
present comparison sits. The narrative is organised around a single recurring question: \emph{how
is word order represented?} Each major advance in the field can be read as a different answer to
that question, and the two architectures compared in this report represent the two most recent
answers.

\section{The sparse feature era}

The earliest statistical approaches to text classification represented a document as a vector of
term counts or weighted term frequencies, discarding word order entirely. The bag-of-words model,
typically weighted by term frequency--inverse document frequency (TF--IDF), paired with a linear
classifier such as multinomial Naive Bayes or a support vector machine, formed the standard
baseline for many years \cite{joachims1998}.

These methods remain surprisingly competitive on topic classification specifically. Topic is
largely signalled by lexical choice: the presence of the words \emph{quarterly}, \emph{earnings}
and \emph{shares} is close to sufficient evidence for a business story, irrespective of the order
in which those words appear. This observation matters for the present study, because it implies
that topic classification is a task where sequential modelling has comparatively little to
contribute --- a point returned to in Chapter~\ref{ch:discussion}.

The principal weaknesses of the sparse approach are the dimensionality of the feature space, the
absence of any notion of semantic similarity between distinct terms, and the inability to
represent negation, modification or any other order-dependent phenomenon.

\section{Distributed word representations}

The second weakness was addressed by distributed representations. Word2Vec \cite{mikolov2013}
and GloVe \cite{pennington2014} learn dense, low-dimensional vectors such that distributionally
similar words occupy nearby regions of the space. A document can then be represented by, for
example, the mean of its word vectors --- a representation that generalises across synonyms in a
way that sparse indicators cannot.

Word order, however, remains unmodelled. Averaging word vectors is still a bag-of-words
operation. The FastText classifier \cite{joulin2017} demonstrated that this simple approach,
with the addition of character n-gram features, achieves strong results on topic classification at
a tiny fraction of the computational cost of neural sequence models --- again suggesting that the
marginal value of sequence modelling on this class of task is modest.

\section{Recurrent neural networks}

Recurrent neural networks introduced genuine sequential modelling. An RNN maintains a hidden state
$h_t$ that is updated at each timestep as a function of the current input and the previous state,
so that in principle $h_T$ summarises the entire sequence.

In practice, the simple recurrent formulation suffers from the vanishing and exploding gradient
problem, analysed by Bengio et al.\ \cite{bengio1994} and later in detail by Pascanu et al.\
\cite{pascanu2013}. Because backpropagation through time multiplies Jacobian matrices once per
timestep, gradient magnitude decays or grows geometrically with the distance over which credit
must be assigned. The practical consequence is that simple RNNs fail to learn dependencies
spanning more than a handful of timesteps. The formal statement of this result is developed in
Section~\ref{sec:vanishing}.

\section{Long Short-Term Memory}

The Long Short-Term Memory architecture of Hochreiter and Schmidhuber \cite{hochreiter1997}
addresses the gradient pathology directly. LSTM introduces a separate \emph{cell state} that is
modified by additive rather than multiplicative interactions, and three multiplicative gates ---
input, forget and output --- that control the flow of information into, within and out of the
cell. Gers et al.\ \cite{gers2000} subsequently added the forget gate in the form now considered
standard.

The critical structural property is that the cell state path from $c_{t-1}$ to $c_t$ passes through
an addition and an elementwise multiplication by the forget gate, rather than through a weight
matrix and a saturating nonlinearity. When the forget gate is near one, gradients propagate
backwards along this path essentially unattenuated. This is the mechanism by which LSTM extends
the effective credit-assignment horizon from a handful of timesteps to hundreds.

The Gated Recurrent Unit \cite{cho2014} simplifies the LSTM to two gates and a single state
vector, with broadly comparable empirical performance and fewer parameters. The present study uses
LSTM rather than GRU because its gate structure is the canonical formulation and its correspondence
with the reference implementation used for verification (Chapter~\ref{ch:verification}) is exact.

\section{Bidirectionality}

A unidirectional recurrent network computes a representation of each position conditioned only on
the preceding context. For classification --- where the entire document is available at inference
time and there is no requirement to generate left to right --- this is an unnecessary restriction.
Bidirectional recurrent networks \cite{schuster1997} run two independent recurrences, one in each
direction, and concatenate the resulting representations, so that each position is contextualised
by both its left and right neighbourhoods.

Bidirectionality is standard in recurrent text classification and is used in the present study.
Note that it doubles the sequential computation, which is relevant to the cost analysis in
Chapter~\ref{ch:results}.

\section{Attention}

Attention was introduced by Bahdanau et al.\ \cite{bahdanau2015} in the context of neural machine
translation, as a mechanism to relieve the bottleneck of compressing an entire source sentence
into a single fixed-length vector. Rather than conditioning on a single summary state, the decoder
computes a weighted combination of all encoder states, with weights determined by a learned
compatibility function between the decoder's current state and each encoder position.

Luong et al.\ \cite{luong2015} simplified the compatibility function to a dot product and
systematised the variants. The essential insight --- that a model can learn \emph{where to look}
rather than being forced to summarise everything into one vector --- proved far more general than
its original application.

\section{The Transformer}

Vaswani et al.\ \cite{vaswani2017} took the decisive step of removing recurrence entirely. The
Transformer computes representations using \emph{self}-attention, in which every position attends
to every other position in the same sequence, combined with position-wise feed-forward layers,
residual connections and layer normalisation \cite{ba2016}.

Two properties of this design are central to the argument of this report.

\keypoint{Constant path length.} In a recurrent network, information must traverse $O(n)$
sequential steps to travel between positions $n$ apart. In self-attention, every pair of positions
is connected directly, so the path length is $O(1)$. This is the structural basis of the
Transformer's advantage on long-range dependency.

\keypoint{Parallelism.} Recurrence is inherently sequential: step $t$ cannot be computed until step
$t-1$ is complete. Self-attention has no such constraint; all positions are computed
simultaneously as a batched matrix product. On parallel hardware this translates directly into
reduced wall-clock time per training step.

Because self-attention is permutation-invariant, positional information must be injected
explicitly. The original work uses fixed sinusoidal encodings; learned positional embeddings are a
common alternative. The pre-norm variant of the residual block, in which layer normalisation is
applied before rather than after each sublayer, was shown by Xiong et al.\ \cite{xiong2020} to
yield better-conditioned gradients at initialisation and to remove the need for a learning-rate
warmup schedule. The present study uses the pre-norm arrangement for precisely this reason.

\section{Pretrained language models}

BERT \cite{devlin2019} demonstrated that a Transformer encoder pretrained on a large unlabelled
corpus with a masked language modelling objective, then fine-tuned on a downstream task,
substantially outperforms task-specific architectures trained from scratch. The paradigm now
dominates the field, with RoBERTa \cite{liu2019} and numerous successors refining the recipe.

It is important to be precise about what this result establishes. The gain attributable to BERT on
a task such as AG News comes overwhelmingly from \emph{pretraining} --- from representations
learned on billions of words of unlabelled text --- rather than from self-attention as an
architectural primitive. Conflating the two is a common error, and the present study is explicitly
designed to avoid it: by training both architectures from scratch on identical data, the
comparison isolates the architectural contribution.

\section{Large language models and the current landscape}

The most recent development is the application of large generative language models to
classification, either zero-shot, few-shot, or via instruction tuning. A 2024 study on text
classification in the age of large language models \cite{mdpi2024} --- identified in the
case-study brief as the related work for this topic --- examines this landscape and reports a
finding directly relevant here: for well-specified classification tasks with adequate labelled
data, compact supervised models remain competitive with, and far cheaper than, large generative
models. The persistent value of small task-specific classifiers is a recurring conclusion in this
literature, and it motivates continued careful study of exactly the architectures compared in this
report.

\section{Prior results on AG News}

AG News was introduced as a classification benchmark by Zhang et al.\ \cite{zhang2015}, who used
it to evaluate character-level convolutional networks. It has since become a standard
four-class topic classification benchmark. Table~\ref{tab:prior} summarises representative
published accuracies.

\begin{table}[ht]
\centering
\caption[Representative published results on AG News]{Representative published test accuracies on
the AG News benchmark. Figures are as reported in the cited work and are not strictly comparable,
since training budgets, preprocessing and the use of pretraining differ. They are presented to
situate the present study rather than to rank methods.}
\label{tab:prior}
\begin{tabular}{lll}
\toprule
\textbf{Approach} & \textbf{Family} & \textbf{Reported accuracy} \\
\midrule
Bag of words + linear classifier \cite{zhang2015}   & Sparse features   & $\approx 88.8\%$ \\
FastText \cite{joulin2017}                          & Embedding average & $\approx 92.5\%$ \\
Character-level CNN \cite{zhang2015}                & Convolutional     & $\approx 90.5\%$ \\
Word-level CNN / LSTM baselines \cite{zhang2015}    & Neural sequence   & $\approx 90$--$92\%$ \\
BERT-base fine-tuned \cite{devlin2019,liu2019}      & Pretrained Transformer & $\approx 94$--$95\%$ \\
\bottomrule
\end{tabular}
\end{table}

Three observations follow from this table and are important context for the results reported later.

First, the spread between a simple bag-of-words classifier and a fine-tuned BERT model is roughly
six percentage points. This is a comparatively narrow band, and it indicates that AG News is a
task on which lexical evidence carries most of the signal.

Second, the models in the middle of this range --- FastText, character CNNs, word-level recurrent
and convolutional baselines --- are clustered tightly around 90--92\%. Architectures of very
different character arrive at very similar accuracy, which is itself suggestive.

Third, the models at the top of the range are distinguished by pretraining, not by architecture.
The present study, which trains from scratch, should therefore be expected to land in the 90--92\%
band, and does.

\section{Research gap}

The literature establishes that Transformers outperform recurrent networks on a wide range of
tasks, and that pretrained Transformers outperform everything. What is less well documented is a
careful, controlled comparison of the two architectures \emph{at matched capacity}, \emph{without
pretraining}, on \emph{short} text, with \emph{statistical testing} of the resulting difference.

Published comparisons frequently differ in parameter count, training budget, preprocessing or
tokenisation between the systems being compared, which makes the attribution of any observed
difference to architecture unsafe. Significance testing of classification differences is uncommon,
despite accuracy differences of under one percentage point being routinely reported as
improvements.

This report addresses that gap for one specific and practically important regime: short-text topic
classification. It holds data, preprocessing, capacity and training budget fixed; it verifies both
implementations numerically before drawing conclusions from them; and it applies a paired
significance test to the resulting predictions.
